\subsection{Entrenamiento de parámetros mediante LoRA}
\label{sec:sft-parametros}
El ajuste completo requiere calcular y almacenar gradientes y estados del optimizador para todos los parámetros del modelo, lo que aumenta el uso de memoria y el costo del entrenamiento. El ajuste parcial reduce esa carga al actualizar solo un subconjunto, aunque la adaptación queda condicionada por las capas o los módulos seleccionados. Frente a estas restricciones, LoRA incorpora matrices entrenables de rango reducido en las proyecciones elegidas y mantiene congelados sus pesos originales~\cite{hu2022lora}.

Como muestra la figura~\ref{fig:sft-parametros}, una misma entrada recorre dos ramas: la transformación original y la actualización de LoRA. Sus salidas se suman, de modo que el modelo puede modificar su respuesta mediante los parámetros añadidos sin actualizar la matriz base.

\begin{figure}[htbp]
\centering
\begingroup
\singlespacing
\definecolor{loratrain}{HTML}{286B9B}
\definecolor{lorafixed}{HTML}{87929A}
\begin{tikzpicture}[>=Stealth,
  block/.style={draw,line width=0.8pt,rounded corners=3pt,align=center,
    minimum height=1.0cm,minimum width=1.8cm,inner sep=6pt,font=\small},
  train/.style={draw=loratrain,fill=loratrain!9},
  fixed/.style={draw=lorafixed,fill=lorafixed!12},
  flow/.style={->,line width=0.85pt,draw=black!65}]
\node[font=\normalsize] (input) at (0,0) {$x$};
\coordinate (split) at (0.7,0);
\node[block,fixed,minimum width=2.5cm] (base) at (3.3,1.25) {$W_0$\\Congelada};
\node[block,train] (a) at (2.0,-1.2) {$A$\\Entrenable};
\node[block,train] (b) at (4.5,-1.2) {$B$\\Entrenable};
\node[block,minimum width=1.0cm,draw=black!45] (scale) at (6.7,-1.2) {$\alpha/r$};
\node[circle,draw=black!60,minimum size=7mm,inner sep=0pt] (sum) at (8.3,1.25) {$+$};
\node[font=\normalsize] (output) at (9.4,1.25) {$y$};
\draw[draw=black!65,line width=0.85pt] (input) -- (split);
\fill[black!65] (split) circle (1.3pt);
\draw[flow,draw=lorafixed] (split) |- (base.west);
\draw[flow,draw=lorafixed] (base) -- (sum);
\draw[flow,draw=loratrain] (split) |- (a.west);
\draw[flow,draw=loratrain] (a) -- (b);
\draw[flow,draw=loratrain] (b) -- (scale);
\draw[flow,draw=loratrain] (scale.east) -| (sum.south);
\draw[flow] (sum) -- (output);
\node[font=\small,align=center] at (4.7,-2.6)
  {$y=W_0x+\dfrac{\alpha}{r}BAx$};
\end{tikzpicture}
\endgroup

\caption{LoRA: pesos base congelados y matrices $A$ y $B$ entrenables.}
\label{fig:sft-parametros}
\end{figure}

Para una matriz $W_0\in\mathbb{R}^{d_{\mathrm{out}}\times d_{\mathrm{in}}}$, la transformación utiliza los pesos efectivos
\begin{equation}
W_{\mathrm{eff}}=W_0+\frac{\alpha}{r}BA,
\qquad A\in\mathbb{R}^{r\times d_{\mathrm{in}}},
\qquad B\in\mathbb{R}^{d_{\mathrm{out}}\times r},
\label{eq:sft-lora}
\end{equation}
donde $r$ determina el rango de la actualización y $\alpha/r$ controla su escala. La matriz $A$ proyecta la entrada a una dimensión reducida y $B$ la lleva a la dimensión de salida, mientras $W_0$ permanece fija.

Durante el SFT, la pérdida se calcula sobre las predicciones obtenidas con esta transformación. A continuación, la retropropagación calcula los gradientes de $A$ y $B$, y el optimizador actualiza esas matrices para reducir la pérdida. Los pesos de $W_0$ siguen interviniendo en el cálculo de la salida, pero no reciben actualizaciones. Este procedimiento se aplica a cada proyección seleccionada para incorporar LoRA.

Para esa matriz, la cantidad de parámetros entrenables añadidos es
\begin{equation}
P_{\mathrm{LoRA}}=r(d_{\mathrm{in}}+d_{\mathrm{out}}),
\qquad P_{W_0}=d_{\mathrm{in}}d_{\mathrm{out}},
\label{eq:sft-lora-parametros}
\end{equation}
por lo que un rango pequeño reduce el número de parámetros entrenables y los recursos necesarios para calcular y almacenar sus gradientes y estados del optimizador, disminuyendo así el costo del entrenamiento. La selección de proyecciones, el rango y la escala definen la configuración de los adaptadores.
